% =====================================================================
%  Protokoll: Modellierung und Simulation eines Photovoltaikmoduls
%  VU Kontinuierliche Simulation, 325.118
%
%  Aufbau: jedes Kapitel liegt in kapitel/ als eigene Datei.
%  Das ist Absicht, damit vier Leute gleichzeitig in Overleaf arbeiten
%  koennen, ohne sich in einer Datei zu blockieren.
%
%  Kompilieren: pdflatex -> bibtex -> pdflatex -> pdflatex
% =====================================================================
\documentclass[11pt,a4paper]{scrartcl}

\usepackage[utf8]{inputenc}
\usepackage[T1]{fontenc}
\usepackage[ngerman]{babel}
\usepackage{lmodern}
\usepackage{microtype}

\usepackage{amsmath,amssymb}
\usepackage{siunitx}
\usepackage{graphicx}
\usepackage{booktabs}
\usepackage{longtable}
\usepackage{caption}
\usepackage{subcaption}
\usepackage{listings}
\usepackage{xcolor}
\usepackage[hidelinks]{hyperref}
\usepackage{cleveref}

\graphicspath{{abbildungen/}{../matlab/figures/}}

% --- siunitx: Einheiten immer ueber \si oder \SI setzen ---------------
\sisetup{
  locale = DE,
  per-mode = symbol,
  range-units = single
}

% --- Einheitliche Symbole, damit alle dasselbe schreiben --------------
\newcommand{\Tm}{T_\mathrm{m}}
\newcommand{\Tamb}{T_\mathrm{amb}}
\newcommand{\Tsky}{T_\mathrm{sky}}
\newcommand{\Tref}{T_\mathrm{ref}}
\newcommand{\Cm}{C_\mathrm{m}}
\newcommand{\Qsolar}{Q_\mathrm{solar}}
\newcommand{\Qkonv}{Q_\mathrm{konv}}
\newcommand{\Qrad}{Q_\mathrm{rad}}
\newcommand{\Wel}{W_\mathrm{el}}

% --- Codelistings -----------------------------------------------------
\lstset{
  language=Matlab,
  basicstyle=\ttfamily\footnotesize,
  keywordstyle=\color{blue!60!black},
  commentstyle=\color{green!40!black},
  breaklines=true,
  frame=single,
  numbers=left,
  numberstyle=\tiny
}

\title{Modellierung und Simulation eines Photovoltaikmoduls}
\subtitle{VU Kontinuierliche Simulation, 325.118}
\author{Vorname Nachname \and Vorname Nachname \and Vorname Nachname \and Vorname Nachname}
\date{\today}

\begin{document}

\maketitle
\tableofcontents
\newpage

\section{Einleitung und Aufgabenstellung}
\label{sec:einleitung}

% VERANTWORTLICH: ...
% WAS HIER HINGEHOERT:
%   - Wozu wird das Modell gebaut? Welche Frage soll es beantworten?
%   - Kurze Nennung der acht Aufgabenpunkte und wo sie im Protokoll stehen.
%   - KEINE Modellgleichungen, die kommen in Abschnitt 2.
% LAENGE: etwa eine Seite.

Die Leistung eines Photovoltaikmoduls haengt neben der Einstrahlung stark von
der Modultemperatur ab. Ziel dieser Arbeit ist ein thermisches Modell, das aus
Globalstrahlung $G$ und Umgebungstemperatur $\Tamb$ die Modultemperatur $\Tm$
und daraus die elektrische Leistung $\Wel$ vorhersagt.

\subsection{Zweck des Modells}
% Diese Frage muss vor allem anderen beantwortet sein. Die Vorlesung
% definiert Modellvalidierung als "erfuellt das Modell seinen Zweck".
% Ohne festgelegten Zweck ist Kapitel 5 nicht schreibbar.

\subsection{Aufbau der Arbeit}

\section{Modellbildung}
\label{sec:modellbildung}

% VERANTWORTLICH: ...
% Das ist das wichtigste Kapitel. Hier steht das "Warum".

\subsection{Systemabgrenzung und Blockschaltbild}
\label{sec:abgrenzung}

% Phase 1 der Modellbildung nach der Vorlesung: Aufteilung in Subsysteme,
% Definition der Systemgrenze, Klaerung der Ein- und Ausgangsgroessen.
% Ergebnis ist das Blockschaltbild aus der Angabe.
% Explizit benennen: Bilanzraum ist das Modul ohne Rahmen.

Eingaenge sind $G$ und $\Tamb$, Ausgaenge sind $\Tm$ und $\Wel$.

\begin{figure}[htbp]
  \centering
  % \includegraphics[width=0.6\textwidth]{blockschaltbild}
  \caption{Whitebox-Modell des Photovoltaikmoduls mit den Eingangsgroessen
           Globalstrahlung und Umgebungstemperatur sowie den Ausgangsgroessen
           Modultemperatur und elektrischer Leistung.}
  \label{fig:blockschaltbild}
\end{figure}

\subsection{Energiebilanz und Wahl des Zustands}
\label{sec:bilanz}

% Phase 2: Naturgesetze anwenden. Hier der erste Hauptsatz.
% Warum genau ein Zustand: das Modul besitzt genau einen unabhaengigen
% Energiespeicher, naemlich seine Waermekapazitaet. Die Ordnung des
% Zustandsvektors folgt damit direkt aus der Physik und ist nicht frei
% gewaehlt. Dieser Absatz ist der direkteste Bezug zur Vorlesung.

\begin{equation}
  \Cm \frac{\mathrm{d}\Tm}{\mathrm{d}t}
    = \Qsolar - \Wel - \Qkonv - \Qrad
  \label{eq:bilanz}
\end{equation}

\subsection{Einstrahlung und elektrische Leistung}
\label{sec:elektrisch}

\begin{equation}
  \Qsolar = \alpha\, G\, A
  \label{eq:qsolar}
\end{equation}

\begin{equation}
  \Wel = G\, A\, (\tau\alpha)\, \eta_\mathrm{ref}
         \left[ 1 - \beta_\mathrm{ref} (\Tm - \Tref) \right]
  \label{eq:wel}
\end{equation}

% Wichtig fuers "Warum": \Wel haengt von \Tm ab. Genau diese Rueckkopplung
% macht aus der Bilanz eine Differentialgleichung. Ohne sie waere \Tm eine
% explizite Formel. Das gehoert ausformuliert.

\subsection{Konvektion}
\label{sec:konvektion}

\begin{equation}
  \Qkonv = h\, A_\mathrm{konv} (\Tm - \Tamb),
  \qquad h = h_a + h_b\, v
  \label{eq:qkonv}
\end{equation}

% Begruendung des linearen Ansatzes gegenueber dem Nusselt-Ansatz von
% Tuncel et al. Argumente: Windrichtung ist in den GeoSphere-Daten nicht
% zuverlaessig nutzbar, die vier Stroemungsregime sind nicht validierbar,
% der Rechenaufwand steht in keinem Verhaeltnis zur erreichbaren Genauigkeit.
% Diese Entscheidung ist der Kern von Aufgabenpunkt 1 ("sinnvolle Annahmen
% und Vereinfachungen") und muss ausformuliert werden.

\subsection{Strahlung}
\label{sec:strahlung}

\begin{equation}
  \Qrad = \sigma A \left[
            \varepsilon_\mathrm{f} (\Tm^4 - \Tsky^4)
          + \varepsilon_\mathrm{b} (\Tm^4 - \Tamb^4) \right],
  \qquad \Tsky = 0{,}0552\, \Tamb^{1{,}5}
  \label{eq:qrad}
\end{equation}

\subsection{Waermekapazitaet}
\label{sec:kapazitaet}

\begin{equation}
  \Cm = A \sum_n d_n\, \rho_n\, c_{p,n}
  \label{eq:cm}
\end{equation}

\subsection{Getroffene Annahmen}
\label{sec:annahmen}

% Kernstueck des Protokolls. Jede Zeile: was wurde angenommen, warum ist
% das zulaessig, was kostet es voraussichtlich an Genauigkeit.
% Die Tabelle waechst waehrend des Projekts mit. Neue Annahme im Code
% heisst neue Zeile hier.

\begin{table}[htbp]
  \centering
  \caption{Getroffene Annahmen und Vereinfachungen mit Begruendung und
           erwarteter Auswirkung auf das Ergebnis.}
  \label{tab:annahmen}
  \begin{tabular}{p{0.28\textwidth} p{0.36\textwidth} p{0.28\textwidth}}
    \toprule
    Annahme & Begruendung & Erwartete Auswirkung \\
    \midrule
    Modul als ein Temperaturknoten (Lumped) &
      Laminatdicke rund \SI{3}{\milli\metre}, Biot-Zahl klein &
      Kein Temperaturgradient abbildbar, im Mittel vernachlaessigbar \\
    Waermeleitung vernachlaessigt &
      folgt aus der Lumped-Annahme &
      gering \\
    Bodentemperatur gleich $\Tamb$ &
      keine Messdaten verfuegbar &
      unterschaetzt Rueckseitenverluste bei aufgeheiztem Untergrund \\
    Lineare Konvektion statt Nusselt-Ansatz &
      siehe \cref{sec:konvektion} &
      \dots \\
    \dots & \dots & \dots \\
    \bottomrule
  \end{tabular}
\end{table}

\section{Modellparameter}
\label{sec:parameter}

% VERANTWORTLICH: ...
% Aufgabenpunkt 3. Die Angabe weist ausdruecklich darauf hin, dass die
% Parameter in den beiden Papers nicht vollstaendig enthalten sind und
% dass auf gute Dokumentation zu achten ist.

\subsection{Vorgehen bei der Recherche}

% Wo wurde gesucht, welche Quellen wurden verworfen und warum.
% Bei Parametern mit stark streuenden Literaturwerten: Bandbreite nennen
% und auf Kapitel 7 verweisen, dort wird die Streuung untersucht.

\subsection{Parameteruebersicht}

% Diese Tabelle muss zeilenweise zu init_parameters.m passen.
% Bei jeder Aenderung im Code hier nachziehen.

\begin{longtable}{l l l p{0.30\textwidth} l}
  \caption{Verwendete Modellparameter mit Quelle und Unsicherheit.}
  \label{tab:parameter} \\
  \toprule
  Symbol & Wert & Einheit & Quelle & Unsicherheit \\
  \midrule
  \endfirsthead
  \toprule
  Symbol & Wert & Einheit & Quelle & Unsicherheit \\
  \midrule
  \endhead
  \bottomrule
  \endfoot
  $A$                     & 1{,}634 & \si{\metre\squared}      & Datenblatt & gering \\
  $\eta_\mathrm{ref}$     & 0{,}127 & --                       & \cite{tuncel2020} & gering \\
  $\beta_\mathrm{ref}$    & 0{,}0045 & \si{\per\kelvin}        & \cite{tuncel2020} & mittel \\
  $\Tref$                 & 298{,}15 & \si{\kelvin}            & STC, IEC 61215 & -- \\
  $\alpha$                & 0{,}90  & --                       & TODO & mittel \\
  $\varepsilon_\mathrm{f}$& 0{,}85  & --                       & TODO & mittel \\
  $h_a$                   & 5{,}7   & \si{\watt\per\metre\squared\per\kelvin} & TODO & hoch \\
  $h_b$                   & 3{,}8   & \si{\watt\per\metre\squared\per\kelvin} & TODO & hoch \\
  $\Cm$                   & TODO    & \si{\joule\per\kelvin}   & aus \cref{eq:cm} & hoch \\
\end{longtable}

\subsection{Herleitung der Waermekapazitaet}

% Schichttabelle mit Dicke, Dichte, spezifischer Waermekapazitaet und Quelle.
% Tuncel et al. weisen darauf hin, dass die Herstellerangaben zu den Schichten
% meist fehlen. Das ist ein guter Aufhaenger fuer Kapitel 7.

\section{Numerische Umsetzung}
\label{sec:numerik}

% VERANTWORTLICH: ...

\subsection{Programmstruktur}

% Kurze Uebersicht der Dateien und wer was aufruft. Ein Absatz plus Liste
% reicht, der Code selbst gehoert nicht ins Protokoll.
% Erwaehnen: Rechnen und Plotten sind getrennt, Ergebnisse werden in
% .mat-Dateien zwischengespeichert.

\subsection{Wahl des Loesungsverfahrens}
\label{sec:solver}

% Das ist der Abschnitt, in dem die Vorlesung am direktesten geprueft wird.
% Zu beantworten:
%   - Ist das Problem steif? Das System ist skalar, ein Steifheitsverhaeltnis
%     ist damit gar nicht definiert. Daraus folgt die Wahl eines expliziten
%     Verfahrens.
%   - Warum ode45 (Dormand-Prince, Ordnung 5 mit eingebetteter Ordnung 4
%     zur Schrittweitensteuerung) und nicht ode15s.
%   - Welche Toleranzen wurden gesetzt und warum nicht die Defaults.
%   - Nachweis der Loesungsunabhaengigkeit: Rechnung mit strengerer Toleranz
%     wiederholen und zeigen, dass sich das Ergebnis nicht mehr aendert.

\subsection{Behandlung der Wetterdaten}

% Warum die Messreihen interpoliert werden muessen: der Solver waehlt seine
% Zeitschritte adaptiv und fragt Zeitpunkte zwischen den Messwerten ab.
% Warum lineare und nicht kubische Interpolation.

\section{Validierung}
\label{sec:validierung}

% VERANTWORTLICH: Linus
% Aufgabenpunkt 4.
% ROHFASSUNG. Abschnitt 5.1 steht, 5.2 und 5.3 warten auf die
% Gruppenentscheidung zur Datengrundlage (siehe doku/annahmen.md,
% offener Punkt 1).

\subsection{Vorgehen und Fehlermasse}

Die Vorlesung nennt drei Kriterien fuer die Validierung eines Modells:
Reproduzierbarkeit, Praediktion und innere Struktur. Auf das vorliegende
Modell angewandt bedeutet das:

\begin{itemize}
  \item \emph{Reproduzierbarkeit.} Bei gleichen Eingangsdaten liefert die
        Implementierung dasselbe Ergebnis. Da alle Parameter an einer
        einzigen Stelle liegen und der Loeser deterministisch arbeitet, ist
        dieses Kriterium konstruktiv erfuellt. Der Nachweis der
        Loesungsunabhaengigkeit in Abschnitt~\ref{sec:solver} gehoert
        ebenfalls hierher.
  \item \emph{Praediktion.} Das Modell muss Messwerte vorhersagen, die
        nicht in seine Aufstellung eingeflossen sind. Dies ist der Kern
        dieses Kapitels.
  \item \emph{Innere Struktur.} Nicht nur das Ergebnis, auch die einzelnen
        Terme muessen physikalisch sinnvoll sein. Zwei Beobachtungen stuetzen
        das: In den Nachtstunden liegt die Modultemperatur um zwei bis drei
        Kelvin \emph{unter} der Lufttemperatur. Das ist kein Rechenfehler,
        sondern die Abstrahlung gegen den nach Swinbank deutlich kaelteren
        Himmel, und es bestaetigt den Strahlungsterm unabhaengig vom
        Gesamtergebnis. Ausserdem schliesst die Energiebilanz ueber eine
        Woche auf ein Hundertstel Prozent, siehe
        Abschnitt~\ref{sec:solver}.
\end{itemize}

\paragraph{Wahl der Fehlermasse.}
Verwendet werden der mittlere absolute Fehler (MAE), die Wurzel des
mittleren quadratischen Fehlers (RMSE) und der mittlere systematische Fehler
(MBE), jeweils in Kelvin. Die drei Masse beantworten verschiedene Fragen:
der MAE gibt die typische Abweichung an, der RMSE gewichtet einzelne grosse
Ausreisser staerker, und der MBE zeigt als einziger die \emph{Richtung} der
Abweichung. Ein Modell kann einen kleinen RMSE aufweisen und trotzdem
systematisch zu tief liegen; erst der MBE macht das sichtbar. Das Residuum
ist als $r = T_\mathrm{sim} - T_\mathrm{mess}$ definiert, ein negativer MBE
bedeutet also eine Unterschaetzung. Diese Konvention entspricht der von
Tuncel et al.~\cite{tuncel2020}, deren Werte damit unmittelbar
vergleichbar sind.

\paragraph{Getrennte Auswertung fuer Tag und Nacht.}
Die Fehlermasse werden zusaetzlich getrennt fuer die Tages- und die
Nachtstunden angegeben. Der Grund dafuer laesst sich am Referenzpaper selbst
ablesen: Tuncel et al. berichten fuer denselben Datensatz einen MAE von
$0{,}90\,\mathrm{K}$ ueber den gesamten Zeitraum, aber $2{,}61\,\mathrm{K}$
allein fuer die Tagesstunden. Nachts liegt das Modul nahe an der
Umgebungstemperatur, die Dynamik ist gering und die Abweichung klein. Wer
ueber volle Tage mittelt, verduennt die Abweichung der Mittagsstunden mit
vielen unauffaelligen Nachtwerten und erhaelt eine Zahl, die guenstiger
aussieht, als das Modell ist. Als Grenze gilt eine Globalstrahlung von
$20\,\mathrm{W/m^2}$; das Referenzpaper definiert die Grenze nicht, die
Festlegung ist daher in der Annahmentabelle dokumentiert.

\paragraph{Kriterium fuer eine bestandene Validierung.}
% TODO In der Gruppe bestaetigen, siehe doku/annahmen.md, offener Punkt 2.
Das Kriterium wurde festgelegt, \emph{bevor} Ergebnisse vorlagen. Wer die
Schwelle erst nach dem Ergebnis waehlt, waehlt die guenstigste. Als Grenze
gilt ein MAE der Tagesstunden von $5\,\mathrm{K}$. Dieser Wert ist aus der
Literatur begruendet: Tuncel et al. erreichen mit dem vollstaendigen
Nusselt-Ansatz und vier unterschiedenen Stroemungsregimen tagsueber einen MAE
von $2{,}61\,\mathrm{K}$. Das hier verwendete Basismodell fasst freie und
erzwungene Konvektion in einem linearen Ansatz zusammen und kann diesen Wert
systematisch nicht unterbieten. Ein Kriterium in der genannten
Groessenordnung ist damit sachlich hergeleitet und nicht nachtraeglich
passend gewaehlt.

\subsection{Datengrundlage}
\label{sec:validierung_daten}

% TODO Dieser Abschnitt haengt an der Gruppenentscheidung, siehe
% doku/annahmen.md, offener Punkt 1. Der folgende Absatz haelt den Befund
% fest, der die urspruengliche Planung hinfaellig gemacht hat.

Als Referenz war Abbildung~1(a) aus Tuncel et al.~\cite{tuncel2020}
vorgesehen, die gemessene und modellierte Modultemperatur ueber fuenf Tage
zeigt. Bei der Vorbereitung der Digitalisierung stellte sich heraus, dass
die Abbildung ausschliesslich die Zielgroesse enthaelt. Globalstrahlung,
Umgebungstemperatur und Windgeschwindigkeit, also genau die drei
Eingangsgroessen des Modells, sind weder dort noch an anderer Stelle der
Arbeit als Zahlenwerte veroeffentlicht; Abschnitt~2.4 der Arbeit haelt nur
fest, dass sie stuendlich an der Wetterstation der Anlage erfasst wurden.
Aus der Abbildung laesst sich damit gewinnen, was das Modell vorhersagen
soll, nicht aber das, womit es rechnet.

\subsection{Ergebnisse}

\begin{figure}[htbp]
  \centering
  % \includegraphics[width=\textwidth]{validierung_temperatur}
  \caption{Simulierte und im Referenzpaper gemessene Modultemperatur ueber
           fuenf Tage. Die Abweichung ist an den Mittagsstunden der klaren
           Tage am groessten.}
  \label{fig:validierung}
\end{figure}

% TODO Zahlen einsetzen, sobald die Datengrundlage geklaert ist.
% Die Konsolenausgabe von plot_validation.m liefert die Werte fuer die
% Bildunterschrift bereits fertig formatiert.

\subsection{Diskussion der Abweichungen}

% TODO Quantifizieren statt entschuldigen: welcher Effekt erklaert wie viel
% Kelvin? Zu behandeln sind mindestens:
%   - fehlende Angaben zu den Eingangsgroessen im Referenzpaper
%   - linearer Konvektionsansatz statt Nusselt-Ansatz
%   - Annahme A_conv = 2A ueber die Einbausituation
%   - ein Temperaturknoten statt Schichtmodell
%   - Ableseungenauigkeit beim Digitalisieren

\section{Anwendungsfall mit Messdaten der GeoSphere Austria}
\label{sec:anwendungsfall}

% VERANTWORTLICH: ...
% Aufgabenpunkt 5, Zeitraum 24.06. bis 01.07.2019.

\subsection{Datengrundlage}

% Station, Zeitaufloesung, Zeitzone, Umgang mit Luecken.
% Falls Horizontalstrahlung verwendet wird und nicht die Strahlung in
% Modulebene: das ist eine Annahme und gehoert in Tabelle 2.1.

\subsection{Ergebnisse}

\begin{figure}[htbp]
  \centering
  % \includegraphics[width=\textwidth]{anwendungsfall_bilanz}
  \caption{Zeitlicher Verlauf der einzelnen Terme der Waermebilanz nach
           \cref{eq:bilanz}. Der konvektive Verlust dominiert waehrend der
           Tagesstunden, der Speicherterm bleibt klein.}
  \label{fig:bilanz}
\end{figure}

\subsection{Dominanter Verlustterm}

% Die Angabe fragt explizit: welcher Verlustterm spielt die groesste Rolle
% und was folgt daraus fuer einen effizienteren Betrieb?
% Antwort quantitativ belegen, ueber die integrierten Energiemengen, nicht
% nur ueber den Blick auf die Kurve.
% Daraus ableiten: Massnahmen zur besseren Kuehlung, Hinterlueftung,
% Montageart, und was das an Mehrertrag bringen wuerde.

\section{Sensitivitaetsanalyse}
\label{sec:sensitivitaet}

% VERANTWORTLICH: ...
% Aufgabenpunkt 6.

\subsection{Vorgehen}

% One-at-a-time: ein Parameter wird variiert, die uebrigen bleiben fest.
% Grenzen dieses Vorgehens benennen: Wechselwirkungen zwischen Parametern
% werden damit nicht erfasst.
% Bewertungsgroessen festlegen, z.B. Maximaltemperatur und Energieertrag.

\subsection{Ergebnisse}

\subsection{Einordnung}

% Guter Diskussionspunkt: Tuncel et al. finden bei Stundenwerten praktisch
% keinen Einfluss der Waermekapazitaet und schliessen daraus, dass eine
% stationaere Rechnung genuegen wuerde. Herteleer et al. argumentieren fuer
% Sekunden- bis Minutenaufloesung ueber die thermische Zeitkonstante
% tau = R_eq * C_eq von einigen hundert Sekunden.
% Mit den 10-Minuten-Daten liegt diese Arbeit dazwischen. Zeigen, auf
% welcher Seite das eigene Ergebnis landet, und daraus ableiten, welche
% Parameter ueberhaupt sorgfaeltig recherchiert werden muessen.

\section{Umsetzung in Simulink und Batterieladung}
\label{sec:simulink}

% VERANTWORTLICH: ...
% Aufgabenpunkt 7.

\subsection{Aufbau des Blockschaltbilds}

% Systematik aus der Vorlesung: pro Zustandsvariable ein Integrator,
% davor ein Summenblock, dann die rechte Seite verdrahten.
% Bei nur einem Zustand ist das ein Integrator.
% Auf algebraische Schleifen achten: die Rueckkopplung von T_m auf W_el
% laeuft ueber den Integrator und ist daher unkritisch.

\subsection{Kopplung mit dem Batteriemodell}

\subsection{Vergleich mit der MATLAB-Implementierung}

% Beide Implementierungen mit denselben Eingangsdaten rechnen und die
% Differenz zeigen. Das ist eine echte Verifikation und deutlich mehr wert
% als zwei separate Kurven.
% Keine Simulink-Scopes ins Protokoll. Daten ins Workspace loggen und in
% MATLAB plotten.

\section{Diskussion und Grenzen des Modells}
\label{sec:diskussion}

% VERANTWORTLICH: ...

\subsection{Gueltigkeitsbereich}

% Die Vorlesung betont, dass der Gueltigkeitsbereich auf den Bereich der
% verfuegbaren Daten beschraenkt ist. Konkret werden:
%   - Welcher Temperaturbereich wurde geprueft?
%   - Welche Windgeschwindigkeiten?
%   - Gilt das Modell fuer Regen, Schnee, Reif? Nein, und warum nicht.

\subsection{Was das Modell nicht kann}

\subsection{Sinnvolle Erweiterungen}

\section{Verwendung von KI-Werkzeugen}
\label{sec:ki}

% VERANTWORTLICH: ...
% Die Nutzung ist in der LVA ausdruecklich erlaubt und muss nicht
% verschleiert werden. Was zaehlt, ist Nachvollziehbarkeit.

% Sinnvolle Gliederung:
%   - Welche Werkzeuge wurden eingesetzt?
%   - Wofuer konkret: Codegeruest, Literatursuche, Textueberarbeitung,
%     Fehlersuche?
%   - Wie wurden die Ergebnisse geprueft? Das ist der entscheidende Punkt.
%     Jede uebernommene Formel wurde gegen die Primaerquelle geprueft,
%     jeder Parameter gegen die zitierte Arbeit.
%   - Was wurde verworfen und warum?


\bibliographystyle{unsrt}
\bibliography{literatur}

\appendix
\section{Anhang}
\label{sec:anhang}

\subsection{Uebersicht der Programmdateien}

\begin{table}[htbp]
  \centering
  \caption{Dateien der MATLAB-Implementierung und ihre Aufgabe.}
  \label{tab:dateien}
  \begin{tabular}{l p{0.6\textwidth}}
    \toprule
    Datei & Aufgabe \\
    \midrule
    \texttt{init\_parameters.m}     & alle Parameter, einzige Stelle mit Zahlenwerten \\
    \texttt{pv\_thermal\_ode.m}     & rechte Seite von \cref{eq:bilanz} \\
    \texttt{calc\_w\_el.m}          & \cref{eq:wel} \\
    \texttt{calc\_h\_conv.m}        & \cref{eq:qkonv} \\
    \texttt{calc\_q\_rad.m}         & \cref{eq:qrad} \\
    \texttt{run\_validation.m}      & Rechnung zu \cref{sec:validierung} \\
    \texttt{run\_usecase.m}         & Rechnung zu \cref{sec:anwendungsfall} \\
    \texttt{run\_sensitivity.m}     & Rechnung zu \cref{sec:sensitivitaet} \\
    \texttt{plot\_*.m}              & Abbildungen, laden nur gespeicherte Ergebnisse \\
    \bottomrule
  \end{tabular}
\end{table}

\subsection{Herleitungen}


\end{document}
